\begin{lemma}[Situation 7.3: unstored theta and stored beta]
Fix a finite loopless multigraph with its individual edge identities, a
palette of $k$ colors, a distinguished edge $e_0$, an input coloring,
an exact terminal tree-sequence $T$, a generating coloring, a
positive-rung ETT, a hierarchy base, a hierarchy counterexample $C$, and
a \noderef{HierarchyMinimumCounterexampleContext} for these data.  Let
$P$ be a \noderef{HierarchyReplacementPresentation}, let
$\alpha,\beta$ be colors, and assume
\[
 2\le j:=\noderef{HierarchyPathNumber}(C.\mathit{hierarchy})
   <\ell:=\noderef{HierarchyFullTailLength}(C).
\]
Assume that every vertex has a missing color under every proper coloring
leaving exactly $e_0$ uncolored.  Let $d$ be a
\noderef{HierarchySituationThreeDisjointData} for
$P,\alpha,\beta$, and put
\[
 \sigma=d.\mathit{switched.coloring},\qquad
 H=d.\mathit{switched.counterexample.hierarchy},\qquad
 q=H.\mathit{dividerCount},
\]
\[
 A=\noderef{HierarchyStarredStage}(H,q),\quad
 B=\noderef{HierarchyFullTailPrefix}(C,j-2),\quad
 Y=\noderef{HierarchyFullTailPrefix}(C,j),
\]
\[
 D=\noderef{HierarchyStoredColorsAt}(H,q),\qquad
 U_j=\{c:\exists r\in\operatorname{Fin}|T|,
 |A|\le r<|A|+j,\ \sigma(T.\mathit{edge}(r))=c\}.
\]
Assume that $d.\theta\notin D$ and that condition (9) takes its stored
alternative:
\[
 \beta\in D\cap\noderef{TreeMissingColors}(\sigma,B),\qquad
 \bigl(d.\mathit{switched.counterexample.good}.\Gamma^q_\beta
        \cup\{\beta\}\bigr)\cap U_j=\varnothing.              \tag{S}
\]
Then these data are impossible.
\end{lemma}
\begin{proof}
Write $R=d.\mathit{switched}$,
$\mathcal E_\sigma=R.\mathit{corresponding.replacementETT}$, and
$g=R.\mathit{counterexample.good}$.  In addition to $A,B,Y$ above put
\[
 X=\noderef{HierarchyFullTailPrefix}(C,j-1),\qquad
 Z=\noderef{HierarchyFullTailPrefix}(C,\ell-1),
\]
and write $y_i$ for the last vertex of the exact prefix
$T(y_i)=\noderef{HierarchyFullTailPrefix}(C,i)$ and $y_\ell$ for the
last vertex of $T$.  The fields of
\noderef{HierarchySituationThreeDisjointData} give
\[
 \beta\in\overline\sigma(B),\qquad
 \beta\in\overline\sigma(y_\ell),\qquad
 d.\theta\in\overline\sigma(y_j),                              \tag{1}
\]
\noderef{IsProperEdgeColoring}, \noderef{IsOneEdgeUncolored} at exactly
$e_0$, and an exact
$\sigma$-\noderef{TashkinovAugmentationSequence} from $X$ to $T$.

Apply \noderef{HierarchyCounterexamplePenultimateElementary} to the
properness, exactly-one-uncoloredness, complete lower-rung package,
fewer-divider induction, rung-local lexicographic minimality, and
$2n+11$ fields of $R.\mathit{minimumContext}$.  In the order just
listed these are all six fields of the stated minimum context, so the
conclusion is
\[
 \noderef{ColoredVertexSetElementary}
   (G,\sigma,\noderef{TreeSequenceVertexSet}(Z)).               \tag{2}
\]
The \noderef{SameGoodHierarchy} field of $R$, including its underlying
equality of the hierarchy stages, identifies the exact terminal tree
and all divider stages with those for $C$.  Expanding
\noderef{HierarchyFullTailLength} and
\noderef{HierarchyFullTailPrefix}, and using $2\le j<\ell$, gives
entrywise \noderef{TreeSequenceInitialSegment} relations
\[
             A\preccurlyeq B\prec X\prec Y\preccurlyeq Z\prec T. \tag{3}
\]
The first relation can be equality when $j=2$.  Every relation in (3)
retains the individual displayed edge identities.

We first rederive equation (10) of the paper.  The vertex $y_j$ lies in
$Z$, while every vertex of $X$ has a different displayed index in $Z$,
by (3) and vertex injectivity in \noderef{TreeSequence}.  Equations
(1)--(2) therefore give
\[
 d.\theta\notin\noderef{TreeMissingColors}(\sigma,X),\qquad
 d.\theta\notin\noderef{TreeMissingColors}(\sigma,A).          \tag{4}
\]
Let $n$ be the terminal rung of $\mathcal E_\sigma$ and put
\[
 D_n^\sigma=\noderef{TashkinovStoredColors}
   (\mathcal E_\sigma.\mathit{series.data},n).
\]
At the last stored index, the definition of
\noderef{HierarchyStoredColorsAt} is the literal equality
\[
                  D=D_n^\sigma\setminus
                     \noderef{TreeMissingColors}(\sigma,A).    \tag{5}
\]
Thus $d.\theta\notin D$, (4), and (5) imply
$d.\theta\notin D_n^\sigma$.

Apply \noderef{HierarchyTailMissingColorExclusion} to the ETT
$\mathcal E_\sigma$, the base $R.\mathit{replacementBase}$, its
hierarchy $H$, the elementary prefix (2), index $j$, and the last
assertion in (1).  An edge of color $d.\theta$ between $A$ and $y_j$
would have to be the theorem's unique exception: $q=0$, the edge is the
actual SE/RE connecting seed at position $|A|$, and
$d.\theta\in D_n^\sigma$.  The last membership contradicts the preceding
paragraph.  Hence
\[
 \sigma(T.\mathit{edge}(r))\ne d.\theta
       \quad\text{whenever }|A|\le r<|A|+j.                    \tag{6}
\]
Equations (4)--(6) are equation (10), including the raw-stored
exclusion which removes the divider-zero seed exception.  Elementarity
(2), (1), and (3) also show that $y_j$ is the first theta-missing
vertex, so
\[
 \noderef{HierarchyFirstMissingPosition}(\sigma,T,d.\theta).
      \mathit{val}=|A|+j.                                     \tag{7}
\]

We next choose the paper's color $\gamma_1$.  By (S), $\beta\in D$;
the `twoColors' clause of \noderef{GoodHierarchy} says that
$g.\Gamma^q_\beta$ has cardinality two.  Choose
\[
                         \gamma_1\in g.\Gamma^q_\beta.         \tag{8}
\]
The `avoidance' clause of the same definition places $\gamma_1$ in
\noderef{HierarchyStageMissingColors}$(H,q)$ and outside
\noderef{HierarchySegmentColors}$(H,q,\beta)$.  At $q>0$ the stage
missing set is the missing-color set of $A=T_{n,q}$; at $q=0$ it is the
missing-color set of the terminal ladder tree $T_n^\sigma$, which is an
exact prefix of $A$.  In either case
\[
              \gamma_1\in\overline\sigma(A).                  \tag{9}
\]
On the other hand, $\beta\in D$ and (5) imply
$\beta\notin\overline\sigma(A)$, so $\gamma_1\ne\beta$.
Moreover (8) and $\beta\in D$ put $\gamma_1$ in the union of all last
Gamma pairs, whereas $d.\mathit{thetaOutsideGamma}$ puts $d.\theta$
outside that union.  Therefore
\[
              \gamma_1\ne\beta,\qquad\gamma_1\ne d.\theta.   \tag{10}
\]

The disjointness in (S), expanded through the exact filter defining
$U_j$, gives simultaneously
\[
 \sigma(T.\mathit{edge}(r))\notin\{\gamma_1,\beta\}
       \quad(|A|\le r<|A|+j).                                 \tag{11}
\]
This interval includes the exact first SE/RE seed when $q=0$; no
tail-only inference is being used to fix that edge.  Put
\[
 i_\gamma=\noderef{HierarchyFirstMissingPosition}
     (\sigma,T,\gamma_1),\qquad
 i_\beta=\noderef{HierarchyFirstMissingPosition}(\sigma,T,\beta).
\]
Equations (9), (S), and (3) give
\[
 i_\gamma.\mathit{val}\le |A|<i_\beta.\mathit{val}\le |B|.  \tag{12}
\]
In particular the first gamma witness precedes the first beta witness.
The exclusion required conditionally by \noderef{ChangeLemma}, with
the ordered pair $(\gamma_1,\beta)$, is the restriction of (11) from
$A$ to $i_\beta$.  Both colors are missing on $Z$ by (9), (S), and
(3), and they are distinct by (10).  For $q>0$, (9) is precisely the
case-1 stage-missing alternative of the Change Lemma; for $q=0$, its
proof above gives gamma missing already on $T_n^\sigma$, precisely the
case-2 alternative.  Passing the six minimum-context fields already
used for (2), along with (2), (10)--(12), to the cited lemma yields
\[
 \noderef{TwoColorComponent}(G,\sigma,\gamma_1,\beta,
       \operatorname{inl}(T.\mathit{vertex}(i_\gamma)))
 =
 \noderef{TwoColorComponent}(G,\sigma,\gamma_1,\beta,
       \operatorname{inl}(T.\mathit{vertex}(i_\beta))).       \tag{13}
\]
This is the formally oriented version of the paper's first component
equality.

Let $Q_0$ be the component in (13), and let $Q_\ell$ be the
$\gamma_1$--$\beta$ component through $y_\ell$.  They are disjoint.
Indeed, \noderef{MissingColorIncidenceEndpoint} gives incidence degree
at most one in \noderef{TwoColorIncidenceGraph} at the gamma-missing vertex
$T.\mathit{vertex}(i_\gamma)$ and at the two beta-missing vertices
$T.\mathit{vertex}(i_\beta)$ and $y_\ell$.  These are three distinct
vertices by (3) and (12): the first lies in $A$, the second lies after
$A$ and no later than $B$, and the third lies after $Z$.  If $Q_0$ and
$Q_\ell$ were equal, \noderef{TwoColorComponentSpanningPath} would span
their common connected component by a simple path, whereas
\noderef{TwoColorIncidenceDegreeAtMostTwo} permits only the two ends of
such a path to have degree at most one.  This contradicts the three
distinct endpoints.  If $Q_\ell$ is a singleton, it cannot contain either
of the two distinct vertices in $Q_0$, so the same conclusion is
immediate.  Thus \noderef{TwoColorComponentsDisjoint} gives
\[
                              Q_\ell\cap Q_0=\varnothing.      \tag{14}
\]
The endpoint lemma also makes $Q_\ell$, including its singleton case, a
\noderef{TwoColorPathComponent}.

Define the first terminal switch in the required orientation by
\[
 \mu_2=\noderef{TwoColorComponentSwap}
       (G,\sigma,\gamma_1,\beta,\operatorname{inl}(y_\ell)).   \tag{15}
\]
It is proper by \noderef{TwoColorComponentSwapProper}; the definition
fixes the value $\mathsf{none}$, so exactly $e_0$ remains uncolored.
Disjointness (14) preserves beta missing at its first witness and gamma
missing at its first witness.  At $y_\ell$, beta was missing by (1), so
the swap makes gamma missing there; if the component is a singleton,
both colors were already missing.  Hence
\[
 \beta\in\overline{\mu_2}(T.\mathit{vertex}(i_\beta)),\qquad
 \gamma_1\in\overline{\mu_2}(T.\mathit{vertex}(i_\gamma))
                  \cap\overline{\mu_2}(y_\ell).              \tag{16}
\]

We verify the data needed to reconstruct the full ETT under $\mu_2$.
First, at any $v\in V(Z)$ outside $Q_\ell$, every incident edge-node is
outside $Q_\ell$, hence every incident color is fixed.  If
$\operatorname{inl}v\in Q_\ell$, (14) makes $v$ distinct from the first
gamma and beta witnesses.  Elementarity (2) then makes both colors
present at $v$; properness makes their incident edge identities unique;
connectedness puts both edge-nodes in $Q_\ell$; and (15) exchanges their
values.  Expanding \noderef{PresentColors} and taking complements gives
\[
 \noderef{MissingColors}(G,\mu_2,v)
   =\noderef{MissingColors}(G,\sigma,v)\qquad(v\in V(Z)).      \tag{17}
\]
This also covers a singleton $Q_\ell$, when only the first alternative
occurs.  Equations (11) and (15) show that every exact displayed edge
from $A$ through $Y$ is fixed.

For each $\eta\in D$, its first missing position is also unchanged.
If $\eta=\beta$, equation (17) through the retained first witness in
$B\subseteq Z$ proves this directly.  The equality $\eta=\gamma_1$ is
impossible because $\eta\in D$ excludes missingness on $A$, while (9)
gives that missingness for gamma.  Every other $\eta$ is globally
unaffected by (15).  Therefore
\[
 \noderef{HierarchyFirstMissingPosition}(\mu_2,T,\eta)
 =\noderef{HierarchyFirstMissingPosition}(\sigma,T,\eta)
                         \qquad(\eta\in D).                   \tag{18}
\]

The literal Cases-One-Two discriminator for the ordered pair
$(\gamma_1,\beta)$ follows from (9): gamma is missing on $T_{n,q}=A$
when $q>0$, and on $T_n^\sigma$ when $q=0$.  Apply
\noderef{StableChangeCasesOneTwo} with the properness and
one-uncoloredness above, the six context fields, (2), the path component
$Q_\ell$, and (14).  Retain its
\noderef{HierarchyWeaklyStable} and \noderef{ModColoring} conclusions.

For completeness we now verify every additional split-tail premise.
Restrict the exact old suffix $X\to T$ first at $Y$ using
\noderef{TashkinovAugmentationSequenceRestrict}; this gives the exact
$\sigma$-suffix $Y\to T$.  Equations (16)--(17) put both affected
colors in \noderef{TreeMissingColors}$(\mu_2,Y)$.

For the source-to-$Y$ route, inspect the actual terminal positive
certificate of $\mathcal E_\sigma$ using \noderef{TashkinovSeries} and
\noderef{TashkinovSeriesStep}.  In PE its
\noderef{TashkinovClosure} starts at $T_n^\sigma$.  In SE or RE it
supplies an individual edge $f_n$, the exact singleton equation
$F_n=\{f_n\}$, the \noderef{TreeSequenceAdjoinEdge} seed
$T_n^\sigma+f_n$, and the route from that seed.  The base and hierarchy
fields from \noderef{HierarchyBase} and \noderef{TashkinovHierarchy},
together with the below-final field of
\noderef{ExtendedTashkinovTree}, put $Y$ on this actual route.  In
positive divider count, \noderef{HierarchyWeaklyStable} contains a
\noderef{TreeColoringInvariant} on $A$, so every exact source-to-$A$
entry is fixed.  At divider zero in PE, weak stability is a
\noderef{CombinedStableColoring}; applying
\noderef{CombinedStableClosurePreserved} fixes every displayed edge of
$T_n^\sigma$ and of the base closure and transports that closure route.
The \noderef{OrderedTreeUnion} in the PE clause of the base identifies
their exact entries in $A$.  At divider zero in SE/RE, the source is the
exact adjoined seed, and (11) fixes its extra edge at position $|A|$;
the \noderef{StableColoring} constituent of weak stability supplies the
required vertexwise control on $T_n^\sigma$.  Combining this source
control with the exact
edge fixation from $A$ through $Y$, and combining (17) on $X$ with
$X\prec Y$ and $|Y|=|X|+1$, all hypotheses of
\noderef{TashkinovAugmentationSequenceColoringTransport} hold.  It gives
the exact $\mu_2$ route
\[
 \begin{cases}
 T_n^\sigma\longrightarrow Y,&\text{PE},\\
 T_n^\sigma+f_n\longrightarrow Y,&\text{SE or RE},
 \end{cases}                                                   \tag{19}
\]
and in the second line it retains the same edge $f_n$, the same seed,
and the literal singleton equality.

It remains only to check the Gamma cutoff.  If an affected color lies
in $g.\Gamma^q_\eta$ for $\eta\in D$, it cannot be beta because
$d.\mathit{betaOutsideGamma}$ excludes beta from the union of all last
Gamma pairs.  It must be gamma.  But gamma lies in
$g.\Gamma^q_\beta$ by (8); the pairwise-disjointness clause of
\noderef{GoodHierarchy}, with $\eta,\beta\in D$, forces $\eta=\beta$.
The beta witness on $B\preccurlyeq Y$ in (S) gives
\[
 \noderef{HierarchyFirstMissingPosition}(\sigma,T,\eta).
      \mathit{val}\le |B|\le |Y|,                             \tag{20}
\]
which is the required cutoff.

Relations (3), (11), (14)--(20), the old suffix $Y\to T$, the literal
Cases-One-Two discriminator, and properness and exact
one-uncoloredness are precisely all premises of
\noderef{StableChangeCasesOneTwoSplitTail}, with split tree $Y$.  The
theorem returns a correspondence certificate
\[
 \mathcal C_1:\noderef{CorrespondingExtendedTashkinovTree}
       (G,k,e_0,\mathcal E_\sigma,\mu_2,T),                    \tag{21}
\]
maximum property for its replacement ETT, a base $\mathcal B_1$, a
hierarchy $H_1$, a good witness $g_1$, and
\[
                    s_1:\noderef{SameGoodHierarchy}(g,g_1).   \tag{22}
\]
By (16), the distinct vertices
$T.\mathit{vertex}(i_\gamma)\in V(A)$ and $y_\ell\notin V(Z)$
both miss gamma under $\mu_2$.  Thus the terminal tree is not elementary.
Together with the positive-rung, maximum, hierarchy, and good fields
returned in (21)--(22), this constructs a
\noderef{HierarchyCounterexample} $C_1$ on $\mathcal B_1$.
Apply \noderef{HierarchyReplacementPresentationTrans} to $R$, $\mu_2$,
$\mathcal C_1$, $\mathcal B_1$, $C_1$, and $s_1$.  It gives an actual
original-relative presentation
\[
 R_1:\noderef{HierarchyReplacementPresentation}
       (G,k,e_0,\mathit{input},T,\phi,\mathit{ett},\mathit{base},C,
        \mathit{ctx}),\qquad R_1.\mathit{coloring}=\mu_2.      \tag{23}
\]
This invocation, rather than a manual identification, composes the
intermediate correspondence with the fixed-original correspondence and
transports the entire minimum context.

We record equation (11) under $R_1$.  The same-good equality (22), which
is retained in (23), identifies the last stored set with $D$ and the
paper-defined pair at beta with $g.\Gamma^q_\beta$.  Equation (17)
preserves the beta witness in $B$, and (11) says that every exact edge
from $A$ through $Y$ is fixed.  Hence
\[
 \beta\in D_1\cap\noderef{TreeMissingColors}(\mu_2,B),\qquad
 (g_1.\Gamma^{q_1}_\beta\cup\{\beta\})
      \cap U_j^{\mu_2}=\varnothing.                            \tag{24}
\]
Here $D_1,q_1,A_1,U_j^{\mu_2}$ are the objects transported to $R_1$;
same-good data identify $q_1,A_1,D_1$ with $q,A,D$.  Thus (24) is the
full transported equation (11), not only its missing-color clause.

The first switch does not involve $d.\theta$, and (17) preserves all
missing sets on $Z$.  Consequently (4), (6), and (7) remain true with
$\mu_2$ in place of $\sigma$.  Correspondence in (21) identifies the
raw stored array, while weak stability identifies the terminal-ladder
missing sets, so the computed raw set $D_n$ is unchanged; hence theta
remains outside it.  This is the full transported equation (10).
Furthermore (16) says gamma remains missing at its hierarchy witness
and at $y_\ell$.  Finally (22), $d.\mathit{thetaOutsideGamma}$, and
$d.\theta\notin D$ give
\[
 d.\theta\notin D_1,\qquad
 d.\theta\notin\bigcup_{\eta\in D_1}g_1.\Gamma^{q_1}_\eta.    \tag{25}
\]

We now apply \noderef{ChangeLemma} for the second time, in the ordered
pair $(\gamma_1,d.\theta)$.  Gamma has its retained first witness in
$A$, whereas the full transported equation (10) says that theta first
becomes missing at $y_j$.  Thus
\[
 i_\gamma^{\mu_2}.\mathit{val}\le |A|
   <\noderef{HierarchyFirstMissingPosition}
       (\mu_2,T,d.\theta).\mathit{val}=|A|+j.                  \tag{26}
\]
The gamma exclusion needed from $A$ to the later witness is (24), since
gamma belongs to beta's reserved pair.  Gamma remains missing on the
last stage, or on $T_n$ when $q=0$, by (17) and (9), so the literal
case-1 or case-2 hypothesis holds.  The two colors are distinct by
(10) and both are missing on $Z$.  Passing these facts and the six
minimum-context fields of $R_1$ to the Change Lemma gives
\[
 \noderef{TwoColorComponent}(G,\mu_2,\gamma_1,d.\theta,
   \operatorname{inl}(T.\mathit{vertex}(i_\gamma^{\mu_2})))
 =
 \noderef{TwoColorComponent}(G,\mu_2,\gamma_1,d.\theta,
   \operatorname{inl}(y_j)).                                  \tag{27}
\]

Let $K_0$ be the common component in (27) and let $K_\ell$ be the
gamma--theta component through $y_\ell$.  The same three-endpoint
argument used for (14) applies to the distinct vertices
$T.\mathit{vertex}(i_\gamma^{\mu_2})$, $y_j$, and $y_\ell$: gamma is
missing at the first and third by (16), theta is missing at the second,
and (3) and (26) make the vertices distinct.  The endpoint, spanning
path, degree-at-most-two, and disjoint-component conclusions of
\noderef{MissingColorIncidenceEndpoint},
\noderef{TwoColorComponentSpanningPath},
\noderef{TwoColorIncidenceDegreeAtMostTwo}, and
\noderef{TwoColorComponentsDisjoint} therefore give
\[
                         K_\ell\cap K_0=\varnothing,           \tag{28}
\]
also when $K_\ell$ is a singleton, and make $K_\ell$ a
\noderef{TwoColorPathComponent}.

Define the second terminal switch by
\[
 \mu_3=\noderef{TwoColorComponentSwap}
       (G,\mu_2,\gamma_1,d.\theta,\operatorname{inl}(y_\ell)). \tag{29}
\]
It is proper and leaves exactly $e_0$ uncolored, exactly as for (15).
Disjointness (28) preserves theta missing at $y_j$.  Gamma is missing at
$y_\ell$ under $\mu_2$, so the switch makes theta missing there; in the
singleton case both colors were already missing.  Hence
\[
                 d.\theta\in\overline{\mu_3}(y_j)
                      \cap\overline{\mu_3}(y_\ell).           \tag{30}
\]

We verify the second reconstruction without suppressing its zero-divider
seed.  The pointwise argument proving (17), now using (28) and the
elementarity of $Z$ under $\mu_2$ supplied by (17), gives
\[
 \noderef{MissingColors}(G,\mu_3,v)
   =\noderef{MissingColors}(G,\mu_2,v)\qquad(v\in V(Z)).       \tag{31}
\]
Equations (24) and the transported equation (10) say that neither gamma
nor theta colors any exact displayed edge from $A$ through $Y$, so
(29) fixes every such edge.  In particular, at $q=0$ in SE/RE it fixes
the actual connecting seed, whose identity and singleton connecting-edge
equation were retained in (19).

The old $\mu_2$ suffix needed at this second step is also explicit.
Restrict $d.\mathit{transportedSuffix}$ to $Y$, and apply
\noderef{TashkinovAugmentationSequenceComponentSwap} to the first switch
(15).  Both beta and gamma are missing on $Y$ under $\mu_2$ by
(16)--(17), so the cited theorem transports the exact $\sigma$ suffix,
with all individual new-edge identities, to a $\mu_2$ suffix $Y\to T$.
Equation (31) and (30) put gamma and theta in
\noderef{TreeMissingColors}$(\mu_3,Y)$.

Apply \noderef{StableChangeCasesOneTwo} to the current presentation
$R_1$ and the ordered pair $(\gamma_1,d.\theta)$.  The literal
discriminator is again gamma's retained stage missingness.  Use the path
$K_\ell$, separation (28), properness and one-uncoloredness, the six
minimum-context fields of $R_1$, and its penultimate elementarity from
\noderef{HierarchyCounterexamplePenultimateElementary}.  The returned
weak stability, together with exact edge fixation through $Y$ and (31),
feeds \noderef{TashkinovAugmentationSequenceColoringTransport} exactly
as in (19).  It produces either the PE route from the current terminal
ladder tree to $Y$, or the SE/RE route from the exact current seed to
$Y$, retaining the actual edge and singleton equation.

For every $\eta\in D_1$, neither affected color equals $\eta$: gamma is
missing on $A$ and hence is outside $D_1$, while theta is outside $D_1$
by (25).  Thus (29) leaves the eta-missing predicate unchanged at every
vertex, so its first missing position is unchanged.  Finally, theta is
outside every last Gamma pair by (25).  If gamma belongs to
$g_1.\Gamma^{q_1}_\eta$, pairwise disjointness and its retained
membership in $g_1.\Gamma^{q_1}_\beta$ force $\eta=\beta$; the beta
first-missing witness in (24) lies already in $B\preccurlyeq Y$.  This
proves the second Gamma cutoff.

All premises of \noderef{StableChangeCasesOneTwoSplitTail} are now
verified for $R_1$, split tree $Y$, old coloring $\mu_2$, new coloring
$\mu_3$, and ordered pair $(\gamma_1,d.\theta)$.  It returns a
correspondence certificate $\mathcal C_2$ relative to
$R_1.\mathit{corresponding.replacementETT}$, maximum property, a base
$\mathcal B_2$, a hierarchy, a good witness $g_2$, and
\[
 s_2:\noderef{SameGoodHierarchy}
        (R_1.\mathit{counterexample.good},g_2).                \tag{32}
\]
Equation (30) makes the exact terminal tree non-elementary under
$\mu_3$, so the returned positive-rung, maximum, hierarchy, and good
data form a \noderef{HierarchyCounterexample} $C_2$ on $\mathcal B_2$.
Apply \noderef{HierarchyReplacementPresentationTrans} to
$R_1,\mu_3,\mathcal C_2,\mathcal B_2,C_2,s_2$.  This gives a second
fixed-original presentation $R_2$ with
\[
                         R_2.\mathit{coloring}=\mu_3.          \tag{33}
\]

The inequalities $2\le j<\ell$, presentation (33), and (30) are
exactly a witness that
\[
 j\in\noderef{HierarchyEligibleCommonMissingIndices}
       (G,k,e_0,\mathit{input},T,\phi,\mathit{ett},
        \mathit{base},C,\mathit{ctx}).                         \tag{34}
\]
Apply \noderef{HierarchySituationThreeEligibleIndexBound} to the fixed
data, the two path bounds, the every-vertex-missing hypothesis, and
(34).  It yields
$j<\noderef{HierarchyPathNumber}(C.\mathit{hierarchy})=j$, a
contradiction.  This proves the claim.
\end{proof}
